\chapter{Ethical AI Governance Frameworks}
\label{ch:frameworks}

This chapter provides a detailed analysis of the three governance frameworks integrated into MEDF and describes the process by which their diverse principles are harmonized into a single, unified evaluation model.

\section{Selected Frameworks}
\label{sec:selected}

\subsection{EU Assessment List for Trustworthy AI (ALTAI)}
\label{sec:altai_detail}

The ALTAI framework, developed by the European Commission's High-Level Expert Group on AI~\cite{aihleg2020altai}, translates the seven requirements of the Ethics Guidelines for Trustworthy AI~\cite{aihleg2019ethics} into a structured self-assessment checklist, as detailed in Section~\ref{sec:altai}. In the MEDF implementation, the framework definition (\texttt{eu\_altai.yaml}) contains 6 criteria entries mapped across the six unified dimensions, with a total of 22 requirements. Intrinsic weights are derived from the number of assessment questions associated with each dimension.

\subsection{NIST AI Risk Management Framework (AI RMF)}
\label{sec:nist_detail}

The NIST AI RMF~\cite{nist2023rmf}, developed by the U.S.\ National Institute of Standards and Technology, takes a process-oriented approach organized around four core functions (Govern, Map, Measure and Manage), as detailed in Section~\ref{sec:nist}. In the MEDF implementation, the framework definition (\texttt{nist\_ai\_rmf.yaml}) contains 6 criteria entries mapped across the six unified dimensions, with a total of 19 subcategories, capturing the risk-management perspective that complements the rights-based focus of ALTAI.

\subsection{Singapore Model AI Governance Framework (MGAF)}
\label{sec:mgaf_detail}

The Singapore Model AI Governance Framework (MGAF)~\cite{pdpc2020mgaf} is structured around two guiding principles and four areas of governance, offering a pragmatic industry-focused perspective, as detailed in Section~\ref{sec:mgaf}. In the MEDF implementation, the framework definition (\texttt{sg\_mgaf.yaml}) contains 6 criteria entries mapped across the six unified dimensions, with a total of 20 principles. The MGAF's emphasis on accountability is reflected in a higher intrinsic weight on the accountability dimension.

\section{Harmonization and the Six Unified Dimensions}
\label{sec:harmonization}

A central contribution of the MEDF design is the harmonization of the three source frameworks into a single, coherent evaluation model based on six Unified Dimensions. These dimensions were derived through a systematic analysis of the overlapping themes and requirements across all three frameworks, informed by the broader AI ethics literature~\cite{floridi2018ai4people,jobin2019landscape}. The six Unified Dimensions are:

\begin{enumerate}[leftmargin=2em]
  \item \textbf{Transparency and Explainability.} This dimension assesses the degree to which the AI system's operations, decisions, and underlying logic can be understood by relevant stakeholders. It encompasses concepts such as traceability, interpretability, and clear communication of the system's capabilities and limitations.

  \item \textbf{Fairness and Non-discrimination.} This dimension evaluates whether the AI system produces equitable outcomes and avoids unjust bias or discrimination against individuals or groups based on protected characteristics such as race, gender, or age.

  \item \textbf{Safety and Robustness.} This dimension measures the technical reliability of the AI system, including its resilience to errors, adversarial attacks, and unexpected inputs. It also considers the system's ability to operate safely within its intended environment and to fail gracefully when encountering conditions outside its design parameters.

  \item \textbf{Privacy and Data Governance.} This dimension assesses the system's compliance with data protection principles, including data minimization, purpose limitation, consent management, and the security of personal data throughout its lifecycle.

  \item \textbf{Human Agency and Oversight.} This dimension evaluates the mechanisms for human oversight, intervention, and ultimate authority over the system's decisions, a principle strongly emphasized in the EU framework.

  \item \textbf{Accountability.} A cornerstone of all three frameworks, this dimension assesses the mechanisms for redress, responsibility, and auditability. It evaluates who is responsible when things go wrong and whether there are clear processes for challenging the system's outcomes.
\end{enumerate}

By mapping the specific clauses, principles, and questions from each of the three source frameworks to these six Unified Dimensions, MEDF creates a comprehensive and cohesive evaluation model. This allows the platform to provide a holistic assessment that is grounded in established international standards, while still offering the flexibility to analyze ethical performance at a granular, dimensional level. The framework definitions, including their mapping to the unified dimensions, are stored as version-controlled YAML files within the repository, ensuring that the foundation of the evaluation logic is itself transparent and auditable.
